\subsection{Defining} \label{s:def}
The first step is to define what \md{} should be extracted or applied to data objects. Search engines can crack many file types and extract not only data from the core file, but also in some cases \md{} embedded in the file's header (e.g. in the case of DICOM images). However, standard data is is made far more useful by applying \emph{semantic meaning} to the content. This step is difficult for a general search engine to do as, by its very nature, semantic meaning is often particular to a user or organization and therefore doesn't lend itself to a generic solution.

Creating an ontology, where semantic meaning can be mapped to data, is referred to as a \emph{Data Dictionary}. Multiple such dictionaries can be created, each with meaning peculiar to a function within an organization. For example, one dictionary may contain a grammar common to a particular vertical, such as the field of healthcare, while other dictionaries will be particular to the organization itself, such as codenames used within a company. The sum of the individual data dictionaries creates the compendium necessary to inform the extraction and application step.
